% ============================================================
%  Chapter 5: Evaluation
% ============================================================
%
% AIM: Present your experimental results clearly and analyse them
%      critically. ~6-8 pages.
%
% CONTENTS:
%   5.1 Experimental Setup (dataset, geometry, noise, hardware, hyperparams)
%   5.2 Baseline Comparison (FBP, TV, SIRT vs paper's baselines)
%   5.3 Quantitative Metrics (PSNR, SSIM table -- your results vs paper)
%   5.4 Qualitative Analysis (reconstruction grids, difference maps, zoom-ins)
%   5.5 Ablation Studies (if time: iterations, denoiser choice, policy viz)
%   5.6 Discussion (how close to paper, what explains gaps, what you'd change)
%
% FIGURES (this chapter should be figure-heavy):
%   - Table 5.1: Main quantitative results (PSNR, SSIM)
%   - Fig 5.1: Reconstruction grid (GT, FBP, TV, TFPnP) for 2-3 images
%   - Fig 5.2: Difference maps
%   - Fig 5.3: PSNR vs noise level plot
%   - Fig 5.4: Learned parameter schedules (sigma_k, mu_k vs iteration)
%   - Fig 5.5: Convergence: PSNR vs iteration for fixed vs learned policy
%
% TIPS:
%   - Every figure must be referenced and discussed in the text
%   - Be honest -- if results don't match the paper, explain why
%   - Quantitative first, then qualitative to support the numbers
% ============================================================

\chapter{Evaluation}
\label{chap:results}

% -----------------------------------------------------------
\section{Experimental Setup}
\label{sec:experimental-setup}
% -----------------------------------------------------------

% TODO: Describe your exact experimental configuration.


% -----------------------------------------------------------
\section{Baseline Comparison}
\label{sec:baseline-comparison}
% -----------------------------------------------------------

% TODO: Report your baseline results (FBP, TV, SIRT).


% -----------------------------------------------------------
\section{Quantitative Metrics}
\label{sec:quantitative-metrics}
% -----------------------------------------------------------

% TODO: Main results table and analysis.

% Template:
% \begin{table}[htbp]
%     \centering
%     \caption{Quantitative comparison. PSNR (dB) / SSIM on test set.}
%     \label{tab:main-results}
%     \begin{tabular}{lcccccc}
%         \toprule
%         & \multicolumn{2}{c}{\textbf{5\% noise}}
%         & \multicolumn{2}{c}{\textbf{7.5\% noise}}
%         & \multicolumn{2}{c}{\textbf{10\% noise}} \\
%         \cmidrule(lr){2-3} \cmidrule(lr){4-5} \cmidrule(lr){6-7}
%         \textbf{Method} & PSNR & SSIM & PSNR & SSIM & PSNR & SSIM \\
%         \midrule
%         FBP     & -- & -- & -- & -- & -- & -- \\
%         TV      & -- & -- & -- & -- & -- & -- \\
%         TFPnP   & -- & -- & -- & -- & -- & -- \\
%         \midrule
%         TFPnP (paper) & 24.50 & -- & 24.23 & -- & 23.72 & -- \\
%         \bottomrule
%     \end{tabular}
% \end{table}


% -----------------------------------------------------------
\section{Qualitative Analysis}
\label{sec:qualitative-analysis}
% -----------------------------------------------------------

% TODO: Reconstruction grids, difference maps, zoom-ins.


% -----------------------------------------------------------
\section{Ablation Studies}
\label{sec:ablations}
% -----------------------------------------------------------

% TODO: If time permits.


% -----------------------------------------------------------
\section{Discussion}
\label{sec:discussion}
% -----------------------------------------------------------

% TODO: Critical analysis.